\section{Discussion}
\label{sec:discussion}

\subsection{What the protocol can and cannot establish}

ShiftBench tests sensitivity to declared transformations.  It cannot prove that
a synthetic operator matches a future deployment, nor can an offline log identify
all counterfactual outcomes.  Exposure interventions are especially limited when
propensities are unknown or support is poor~\cite{Schnabel2016,Saito2020}.  We
therefore recommend treating stress tests as evidence about failure modes, not as
a substitute for controlled online experiments.

The protocol also does not prescribe a single scalar leaderboard.  Aggregating
utility, calibration, equity, and coverage requires normative weights that vary
by application.  A Pareto report keeps these trade-offs explicit.  If a scalar
is operationally necessary, weights should be chosen before viewing test results
and accompanied by the disaggregated metrics.

\subsection{Threats to validity}

\paragraph{Construct validity.}
Topic rotation is only one representation of preference drift, and popularity
reweighting may not reproduce a real campaign or news cycle.  Operators should be
calibrated against historical incidents when such data exist.  The diagnostic
then reports the distance between the synthetic shift and the incident.

\paragraph{Internal validity.}
Randomness in splitting, negative sampling, hyperparameter search, and model
training can exceed differences between systems.  Frozen manifests, equal search
budgets, multiple seeds, and paired intervals reduce this risk.  They do not
correct implementation bugs; unit tests should confirm the identity condition
and all declared invariants.

\paragraph{External validity.}
The dry run is intentionally synthetic and does not support claims about any
domain.  Empirical use should span datasets with different feedback semantics,
catalog dynamics, and user populations.  Results from one logging policy should
not be generalized to another without a support analysis.

\paragraph{Metric validity.}
NDCG reflects rank-sensitive relevance but not satisfaction, long-term welfare,
or strategic response.  Calibration can conflict with discovery, and coverage
does not ensure meaningful exposure.  Application-specific outcomes should be
added rather than inferred from these proxies.

\subsection{Responsible use}

Stress tests can expose disparate losses across activity groups, but activity is
not a substitute for protected characteristics.  Where sensitive attributes can
be collected lawfully and ethically, evaluation should be designed with affected
stakeholders and appropriate privacy protections.  The protocol must not be used
to infer sensitive traits from behavior merely to populate a fairness table.
Moreover, a model that performs well in ShiftBench may still create harmful
feedback loops after deployment.
